\documentclass[conference]{IEEEtran}

\usepackage{cite}
\usepackage{amsmath,amssymb,amsfonts}
\usepackage{graphicx}
\usepackage{booktabs}
\usepackage{array}
\usepackage{url}
\usepackage{xcolor}

\title{Notes on Dynamic Programming}

\author{
\IEEEauthorblockN{Luis Mendez}
}

\begin{document}
\maketitle


\section*{Dynamic Programming}

Dynamic programming is a method for solving complex problems by breaking them down into simpler subproblems. In the context of MDPs, dynamic programming algorithms are used to compute the value functions and optimal policies. The two main dynamic programming algorithms for MDPs are:
- Value Iteration: This algorithm iteratively updates the value function until it converges to the
optimal value function. The update rule is given by:
$$ V_{k+1}(s) = \max_{a} \sum_{s'} P(s'|s,a) \left[ R(s,a,s') + \gamma V_k(s') \right] $$
Where $V_k(s)$ is the value function at iteration $k$, and the maximum is taken over all possible actions $a$.
- Policy Iteration: This algorithm alternates between policy evaluation and policy improvement steps. In the policy evaluation step, the value function for the current policy is computed using the Bellman equation. In the policy improvement step, the policy is updated to be greedy with respect to the current value function. The update rule for the policy improvement step is given by:
$$ \pi_{k+1}(s) = \arg\max_{a} \sum_{s'} P(s'|s,a) \left[ R(s,a,s') + \gamma V^{\pi_k}(s') \right] $$
Where $\pi_k(s)$ is the policy at iteration $k$, and the maximum is taken over all possible actions $a$.

\section*{Optimal substructure}
We are looking for the optimal policy for each state, and we can use the optimal policy for the successor states to compute the optimal policy for the current state. This is known as the principle of optimality, which states that an optimal policy has the property that whatever the initial state and initial decision are, the remaining decisions must constitute an optimal policy with regard to the state resulting from the first decision.
The optimal solution to all subproblems can be used to construct the optimal solution to the overall problem. This is a key property that allows dynamic programming algorithms to efficiently solve MDPs.


\begin{figure}[ht]
    \centering
    \includegraphics[width=0.8\linewidth]{image.png}
\end{figure}

Subproblems are not independent, and the optimal solution to one subproblem may depend on the optimal solution to another subproblem. This is why we need to solve all subproblems in order to find the optimal solution to the overall problem.

\subsection*{Overlapping subproblems}
The optimal solution to a subproblem may be used multiple times in the computation of the optimal solution
to the overall problem. This is known as overlapping subproblems, and it is a key property that allows dynamic programming algorithms to efficiently solve MDPs.

\subsection*{Model of the environment}
In order to apply dynamic programming algorithms to solve MDPs, we need to have a model of the environment. This model includes the transition probabilities $P(s'|s,a)$ and the reward function $R(s,a,s')$. The transition probabilities describe how the environment changes in response to the agent's actions, while the reward function describes the immediate reward received by the agent for taking a particular action in a particular state.
Having a model of the environment allows us to compute the value functions and optimal policies using dynamic programming
algorithms. However, in many real-world scenarios, we may not have access to a complete model of the environment, which is why we need to use model-free reinforcement learning algorithms that do not require a model of the environment.

\end{document}